\documentclass[reqno]{amsart}
\usepackage[T1]{fontenc}
\usepackage{amsmath,amssymb,amsthm}
\usepackage{xr-hyper}
\usepackage[hidelinks]{hyperref}
\externaldocument[zr-]{20261125T001000--zeros-to-ramp-sums__workbook}
\externaldocument[zi-]{20261125T001100--inert-primes-and-estimate-b__workbook}
\externaldocument[rc-]{20261125T001500--real-characters-over-prime-moduli__workbook}
\externaldocument[gf-]{20261125T005000--galois-fourier-and-averaging__workbook}
\externaldocument[cs-]{20261125T005500--columns-from-a-sphere__workbook}
\externaldocument[as-]{20261125T021500--assembly-at-one-scale__workbook}

% Chapter number; \ref{chap} gives it.
\newcounter{chap}
% Chapter numbers, generated from 20261125T012000--grh-workbook__guide.org
% by scripts/make-chapters.py.  Do not edit: change the order of the
% index and rerun the script (./preview-site.sh and CI do this).
%
% Usage in a sheet, after \newcounter{chap}:
%   % Chapter numbers, generated from 20261125T012000--grh-workbook__guide.org
% by scripts/make-chapters.py.  Do not edit: change the order of the
% index and rerun the script (./preview-site.sh and CI do this).
%
% Usage in a sheet, after \newcounter{chap}:
%   % Chapter numbers, generated from 20261125T012000--grh-workbook__guide.org
% by scripts/make-chapters.py.  Do not edit: change the order of the
% index and rerun the script (./preview-site.sh and CI do this).
%
% Usage in a sheet, after \newcounter{chap}:
%   \input{chapters}\chapnumber{<sheet ID>}
\makeatletter
\providecommand\chapters@set[2]{\@namedef{chapters@#1}{#2}}
\chapters@set{20261125T001000}{1}% zeros-to-ramp-sums
\chapters@set{20261125T001100}{2}% inert-primes-and-estimate-b
\chapters@set{20261125T005000}{3}% galois-fourier-and-averaging
\chapters@set{20261125T005500}{4}% columns-from-a-sphere
\chapters@set{20261125T000000}{5}% slope-method
\chapters@set{20261126T030100}{6}% a-twist-for-each-prime
\chapters@set{20261125T003000}{7}% the-chain-and-the-square-root
\chapters@set{20261125T010000}{8}% rank-lemma-and-zero-estimate
\chapters@set{20261125T010500}{9}% poincare-series-and-near-orthogonality
\chapters@set{20261124T223000}{10}% elliptic-curves-for-the-proof
\chapters@set{20261125T009500}{11}% fock-model-and-theta-intertwiners
\chapters@set{20261124T223100}{12}% complex-multiplication-and-inert-reduction
\chapters@set{20261125T008000}{13}% one-cm-determinant
\chapters@set{20261124T214500}{14}% supersingular-theta-divisibility
\chapters@set{20261124T214600}{15}% the-isotropic-plane
\chapters@set{20261125T013700}{16}% pairings-on-a-formal-group-of-height-two
\chapters@set{20261125T013500}{17}% hecke-nondegeneracy-on-the-tree
\chapters@set{20261126T030200}{18}% the-bound-is-attained
\chapters@set{20261125T009000}{19}% kronecker-limit-and-chowla-selberg
\chapters@set{20261125T002000}{20}% definite-theta-lifts
\chapters@set{20261125T015500}{21}% siegel-weil-in-the-smallest-case
\chapters@set{20261125T002100}{22}% definite-theta-lifts-and-norms
\chapters@set{20261125T016500}{23}% hecke-operators-and-theta-lifts
\chapters@set{20261125T016700}{24}% divisibility-along-the-chain
\chapters@set{20261125T004000}{25}% denominators-and-sha
\chapters@set{20261125T004100}{26}% tate-twists-and-sha
\chapters@set{20261126T020430}{27}% the-minus-part-by-gauss-sums
\chapters@set{20261125T014100}{28}% congruence-modules
\chapters@set{20261125T014500}{29}% eisenstein-congruences-and-selmer-classes
\chapters@set{20261125T014700}{30}% pseudocharacters-and-the-length-bound
\chapters@set{20261125T011000}{31}% galois-representations-black-box
\chapters@set{20261126T030300}{32}% the-data
\chapters@set{20261126T020210}{33}% bessel-poincare-series-and-their-gram-operators
\chapters@set{20261126T030400}{34}% the-seesaw-identity
\chapters@set{20261126T030500}{35}% divisibility-of-the-composite-lift
\chapters@set{20261126T030600}{36}% theorem-a-short-forms
\chapters@set{20261126T030700}{37}% the-determinant-of-a-chain
\chapters@set{20261126T030800}{38}% the-cost-of-one-lift
\chapters@set{20261126T030900}{39}% denominators-of-projections-to-lifts
\chapters@set{20261126T031000}{40}% twists-with-few-congruences
\chapters@set{20261126T031100}{41}% theorem-c-and-the-main-inequality
\providecommand\chapnumber[1]{%
  \@ifundefined{chapters@#1}%
    {\PackageError{chapters}{Sheet #1 is not listed in the index}%
      {List it in 20261125T012000--grh-workbook__guide.org\MessageBreak
       and run scripts/make-chapters.py.}}%
    {\edef\thechap{\@nameuse{chapters@#1}}}}
\makeatother
\chapnumber{<sheet ID>}
\makeatletter
\providecommand\chapters@set[2]{\@namedef{chapters@#1}{#2}}
\chapters@set{20261125T001000}{1}% zeros-to-ramp-sums
\chapters@set{20261125T001100}{2}% inert-primes-and-estimate-b
\chapters@set{20261125T005000}{3}% galois-fourier-and-averaging
\chapters@set{20261125T005500}{4}% columns-from-a-sphere
\chapters@set{20261125T000000}{5}% slope-method
\chapters@set{20261126T030100}{6}% a-twist-for-each-prime
\chapters@set{20261125T003000}{7}% the-chain-and-the-square-root
\chapters@set{20261125T010000}{8}% rank-lemma-and-zero-estimate
\chapters@set{20261125T010500}{9}% poincare-series-and-near-orthogonality
\chapters@set{20261124T223000}{10}% elliptic-curves-for-the-proof
\chapters@set{20261125T009500}{11}% fock-model-and-theta-intertwiners
\chapters@set{20261124T223100}{12}% complex-multiplication-and-inert-reduction
\chapters@set{20261125T008000}{13}% one-cm-determinant
\chapters@set{20261124T214500}{14}% supersingular-theta-divisibility
\chapters@set{20261124T214600}{15}% the-isotropic-plane
\chapters@set{20261125T013700}{16}% pairings-on-a-formal-group-of-height-two
\chapters@set{20261125T013500}{17}% hecke-nondegeneracy-on-the-tree
\chapters@set{20261126T030200}{18}% the-bound-is-attained
\chapters@set{20261125T009000}{19}% kronecker-limit-and-chowla-selberg
\chapters@set{20261125T002000}{20}% definite-theta-lifts
\chapters@set{20261125T015500}{21}% siegel-weil-in-the-smallest-case
\chapters@set{20261125T002100}{22}% definite-theta-lifts-and-norms
\chapters@set{20261125T016500}{23}% hecke-operators-and-theta-lifts
\chapters@set{20261125T016700}{24}% divisibility-along-the-chain
\chapters@set{20261125T004000}{25}% denominators-and-sha
\chapters@set{20261125T004100}{26}% tate-twists-and-sha
\chapters@set{20261126T020430}{27}% the-minus-part-by-gauss-sums
\chapters@set{20261125T014100}{28}% congruence-modules
\chapters@set{20261125T014500}{29}% eisenstein-congruences-and-selmer-classes
\chapters@set{20261125T014700}{30}% pseudocharacters-and-the-length-bound
\chapters@set{20261125T011000}{31}% galois-representations-black-box
\chapters@set{20261126T030300}{32}% the-data
\chapters@set{20261126T020210}{33}% bessel-poincare-series-and-their-gram-operators
\chapters@set{20261126T030400}{34}% the-seesaw-identity
\chapters@set{20261126T030500}{35}% divisibility-of-the-composite-lift
\chapters@set{20261126T030600}{36}% theorem-a-short-forms
\chapters@set{20261126T030700}{37}% the-determinant-of-a-chain
\chapters@set{20261126T030800}{38}% the-cost-of-one-lift
\chapters@set{20261126T030900}{39}% denominators-of-projections-to-lifts
\chapters@set{20261126T031000}{40}% twists-with-few-congruences
\chapters@set{20261126T031100}{41}% theorem-c-and-the-main-inequality
\providecommand\chapnumber[1]{%
  \@ifundefined{chapters@#1}%
    {\PackageError{chapters}{Sheet #1 is not listed in the index}%
      {List it in 20261125T012000--grh-workbook__guide.org\MessageBreak
       and run scripts/make-chapters.py.}}%
    {\edef\thechap{\@nameuse{chapters@#1}}}}
\makeatother
\chapnumber{<sheet ID>}
\makeatletter
\providecommand\chapters@set[2]{\@namedef{chapters@#1}{#2}}
\chapters@set{20261125T001000}{1}% zeros-to-ramp-sums
\chapters@set{20261125T001100}{2}% inert-primes-and-estimate-b
\chapters@set{20261125T005000}{3}% galois-fourier-and-averaging
\chapters@set{20261125T005500}{4}% columns-from-a-sphere
\chapters@set{20261125T000000}{5}% slope-method
\chapters@set{20261126T030100}{6}% a-twist-for-each-prime
\chapters@set{20261125T003000}{7}% the-chain-and-the-square-root
\chapters@set{20261125T010000}{8}% rank-lemma-and-zero-estimate
\chapters@set{20261125T010500}{9}% poincare-series-and-near-orthogonality
\chapters@set{20261124T223000}{10}% elliptic-curves-for-the-proof
\chapters@set{20261125T009500}{11}% fock-model-and-theta-intertwiners
\chapters@set{20261124T223100}{12}% complex-multiplication-and-inert-reduction
\chapters@set{20261125T008000}{13}% one-cm-determinant
\chapters@set{20261124T214500}{14}% supersingular-theta-divisibility
\chapters@set{20261124T214600}{15}% the-isotropic-plane
\chapters@set{20261125T013700}{16}% pairings-on-a-formal-group-of-height-two
\chapters@set{20261125T013500}{17}% hecke-nondegeneracy-on-the-tree
\chapters@set{20261126T030200}{18}% the-bound-is-attained
\chapters@set{20261125T009000}{19}% kronecker-limit-and-chowla-selberg
\chapters@set{20261125T002000}{20}% definite-theta-lifts
\chapters@set{20261125T015500}{21}% siegel-weil-in-the-smallest-case
\chapters@set{20261125T002100}{22}% definite-theta-lifts-and-norms
\chapters@set{20261125T016500}{23}% hecke-operators-and-theta-lifts
\chapters@set{20261125T016700}{24}% divisibility-along-the-chain
\chapters@set{20261125T004000}{25}% denominators-and-sha
\chapters@set{20261125T004100}{26}% tate-twists-and-sha
\chapters@set{20261126T020430}{27}% the-minus-part-by-gauss-sums
\chapters@set{20261125T014100}{28}% congruence-modules
\chapters@set{20261125T014500}{29}% eisenstein-congruences-and-selmer-classes
\chapters@set{20261125T014700}{30}% pseudocharacters-and-the-length-bound
\chapters@set{20261125T011000}{31}% galois-representations-black-box
\chapters@set{20261126T030300}{32}% the-data
\chapters@set{20261126T020210}{33}% bessel-poincare-series-and-their-gram-operators
\chapters@set{20261126T030400}{34}% the-seesaw-identity
\chapters@set{20261126T030500}{35}% divisibility-of-the-composite-lift
\chapters@set{20261126T030600}{36}% theorem-a-short-forms
\chapters@set{20261126T030700}{37}% the-determinant-of-a-chain
\chapters@set{20261126T030800}{38}% the-cost-of-one-lift
\chapters@set{20261126T030900}{39}% denominators-of-projections-to-lifts
\chapters@set{20261126T031000}{40}% twists-with-few-congruences
\chapters@set{20261126T031100}{41}% theorem-c-and-the-main-inequality
\providecommand\chapnumber[1]{%
  \@ifundefined{chapters@#1}%
    {\PackageError{chapters}{Sheet #1 is not listed in the index}%
      {List it in 20261125T012000--grh-workbook__guide.org\MessageBreak
       and run scripts/make-chapters.py.}}%
    {\edef\thechap{\@nameuse{chapters@#1}}}}
\makeatother
\chapnumber{20261125T020000}
\renewcommand{\thesection}{\thechap.\arabic{section}}

\theoremstyle{definition}
\newtheorem{problem}{Problem}[section]
\theoremstyle{plain}
\newtheorem{theorem}{Theorem}[section]

\newcommand{\Z}{\mathbb{Z}}
\newcommand{\Q}{\mathbb{Q}}
\newcommand{\C}{\mathbb{C}}
\newcommand{\Bc}{\mathcal{B}}
\newcommand{\rd}{\operatorname{rd}}
\newcommand{\De}{\operatorname{De}}
\newcommand{\ol}{\overline}
\renewcommand{\Re}{\operatorname{Re}}

\setlength{\emergencystretch}{2em}

\title{The analytic frame: windows, translates and auxiliary characters}
\author{}
\date{November 25, 2026}

\begin{document}
\maketitle
\refstepcounter{chap}\label{chap}

\paragraph{Prerequisites.}
From Chapter~\ref{zr-chap}: \(G\), \(G_\chi\), \(c_\chi\) and \(G^\psi\) (equations~\eqref{zr-eq:G}, \eqref{zr-eq:Gchi} and Section~\ref{zr-sec:dirichlet}), and Theorems~\ref{zr-thm:intervals} and~\ref{zr-thm:dirichlet-intervals}.  From Chapter~\ref{zi-chap}: Problem~\ref{zi-pr:Rqd} and Theorem~\ref{zi-thm:good-q-interval}.  From Chapter~\ref{rc-chap}: Theorems~\ref{rc-thm:estB-q} and~\ref{rc-thm:linnik}.  From Chapter~\ref{gf-chap}: all of it, in particular \(\Gamma\), \(\Phi\), the characters \(\theta_{\xi,j,l}\), \(P_{t,w}\), \(M_w\), \(R_{q,d}\), \(R_\xi\), \(\mathcal{K}(x)\), and Theorems~\ref{gf-thm:balance} and~\ref{gf-thm:two-scales}.  From Chapter~\ref{cs-chap}: Theorem~\ref{cs-thm:behrend} and Section~\ref{cs-sec:tops}.  The zero-free region, the exceptional zero and the prime number theorem for characters \cite[Chs.~14 and~20]{Dav}; the conductor--discriminant formula \cite[Ch.~3]{Wa}. % check-numbers: ok (Davenport): Chs.~14 and~20

\medskip
The earlier chapters of Part~I proved each analytic step in its simplest setting: one window theorem for \(\zeta\), one field \(E\) with \(H = (\Z/4n\Z)^\times/\{\pm1\}\), one auxiliary prime at a time.  In what generality does the proof use them?  This chapter states the analytic part of the argument as Adriazola proves it \cite[\S\S2--6]{Adr}, as cited theorems, each with its simplest incarnation in the earlier chapters and its status.  The problems specialize the statements, check their exponents, and deduce the forms used later.  Every statement here is proved in \cite{Adr} from published theorems, with one exception: Theorem~\ref{thm:data}(ii) uses the Selmer bound on windows, \cite[Hyp.~17.11]{Adr}, assumed there.  The Lean formalization \cite{Lean} proves the final theorem; the status of each input is recorded in Section~\ref{as-sec:map} of Chapter~\ref{as-chap}. % check-numbers: ok (Adriazola): \S\S2--6

Throughout, \(\psi\) is a primitive Dirichlet character of conductor \(M\) and order \(o\), \(N_0 := 8M\), and \(0 < \varepsilon < \frac1{52}\).  We write \(G_\chi\) for the complex function~\eqref{zr-eq:Gchi}; it is the function \(G^\ast_\chi\) of \cite[(5.1)]{Adr}, whose real part \cite{Adr} calls \(G_\chi\).  We keep the notation of Chapter~\ref{gf-chap}, with \(u\) for the translates: \(\Gamma = \operatorname{Gal}(E/\Q)\), \(\Phi = \operatorname{Gal}(E/K)\) with \(K = \Q(\sqrt{-d})\), \(n = \lvert H \rvert\), \(A_0 \subset H\) of index \(o\), and \(G_\theta = G_{\chi_\theta}\).  The element \(\iota_d \in \Gamma\) acts by \(-1\) on \(\sqrt{-d}\) and trivially on \(\sqrt{-q}\) and on \(H\); it generates the inertia group at \(d\).

\section{Windows}\label{sec:windows}

Chapter~\ref{zr-chap} excluded zeros by an upper bound on windows \([Y, Y^{3/2}]\), via Kaczorowski--Pintz and a separate argument for a real zero.  Adriazola proves one theorem that does both.

\begin{theorem}[{\cite[Thm.~2.4]{Adr}}]\label{thm:window}
Let \(F \colon [1,\infty) \to \mathbb{R}\) be measurable with \(\lvert F(x) \rvert \le Bx\), and suppose \(A(s) := \int_1^\infty F(x)x^{-s-1}\,dx\) extends to a meromorphic function on \(\C\).  Let \(0 \le \theta < 1\).  Suppose that every pole \(p\) of \(A\) with \(\Re p > \theta\) is simple with \(\Re p < 1\), and that the residue at every real pole \(p > \theta\) is positive.  If there are \(\lambda > 1\) and \(C\) such that \(F(x) \le Cx^\theta\) on \([Y, Y^\lambda]\) for arbitrarily large \(Y\), then \(A\) has no pole in \(\Re s > \theta\).
\end{theorem}

\noindent\emph{Simplest incarnation}: Theorem~\ref{zr-thm:intervals}.  \emph{Status}: proved in \cite{Adr}; a quantitative form of Landau's theorem, with Taylor coefficients at a real point as in Section~\ref{zr-sec:real-pole}.  For \(G_+ := G + \sum_i w_i\Re G_{\chi_i}\) with \(w_i > 0\) and nontrivial \(\chi_i\), it gives \cite[Cor.~2.5]{Adr}: a bound \(G_+(x) \le Cx^\theta\) on such windows implies \(\zeta(s) \ne 0\) and \(L(s,\chi_i) \ne 0\) for \(\Re s > \theta\).

\begin{problem}\label{pr:windows}
(i) Theorem~\ref{thm:window} applies to \(F = G\) with \(\lambda = \frac32\), \(C = 1\), and gives Theorem~\ref{zr-thm:intervals}; the residue at a real zero \(\beta \in (\theta, 1)\) of \(\zeta\) of multiplicity \(m\) is \(m/(\beta(1-\beta)) > 0\).
(ii) The corollary with \(\chi_j = \psi_j\) and \(w_j = 1/(3o+1)\) for \(1 \le j < o\) is Theorem~\ref{zr-thm:dirichlet-intervals}.
\end{problem}

\section{The excess and the average}\label{sec:average}

Let \(E\) be the field of Theorem~\ref{thm:field} below, and identify \(\Gamma\) with \(H \times \{\pm1\}^2\) as there.  Let \(c' := ([1], -1, -1)\); it is the complex conjugation \(c\) when \(\psi(-1) = 1\).  Define the cells, \(\Lambda_d\), \(P_{t,w}\) and \(M_w\) as in Chapter~\ref{gf-chap}, with \(c'\) in place of \(c\).  For \(a > 0\), \(t \in H\) and \(w \colon \Phi \to (0,\infty)\), Adriazola's \emph{excess} is
\[
  \De(a, t, w) := \frac{1}{2n}\sum_{g \notin \Phi}\sum_{\theta}\ol{\theta(g)}\, w(t g c')\, G_\theta\Bigl(\frac{a}{w(tgc')}\Bigr) .
\]

\begin{theorem}[excess identity, {\cite[Prop.~3.3 and Cor.~3.4]{Adr}}]\label{thm:excess}
For all \(a > 0\), \(t \in H\) and \(w \colon \Phi \to (0,\infty)\), \(P_{t,w}(a) = M_w(a) + \De(a,t,w)\).  Here \(q, d\) are as in Theorem~\ref{thm:second}, so that \(d \equiv 7 \pmod 8\) and \(d \equiv -1\) modulo \(q\) and modulo every odd prime dividing \(M\); then the primes dividing \(2qM\) split in \(K\) and contribute nothing to \(P_{t,w}\).
\end{theorem}

\noindent\emph{Simplest incarnation}: Theorem~\ref{gf-thm:balance}.  \emph{Status}: proved in \cite{Adr}; an identity between finite sums.

The average removes the characters that are not trivial on \(A_0\) and keeps the \(4o\) characters \(\theta_{\xi,j,l}\), \(\xi \in \langle\psi\rangle\) (Problem~\ref{gf-pr:A0}).  Let \(T_2 \subset A_0\), \(\delta := \lvert T_2\rvert/n\), and call a cell thick for \(u \in A_0\) if it lies in row \(1\) and in a column of \(uT_2\), as in~\eqref{gf-eq:thick-thin}.  Let \(T \subset H\) have \(N_T\) elements, and let \(\kappa, \nu \ge 0\) with \(\lvert\sum_{t\in T}\xi(t) - \kappa N_T\rvert \le \nu N_T\) for \(\xi \in \langle\psi\rangle \smallsetminus 1\); put \(\kappa_1 := 1\) and \(\kappa_\xi := \kappa\) otherwise.

\begin{theorem}[{\cite[Lemma~4.1 and Thm.~4.2]{Adr}}]\label{thm:main-step}
Let \(x \ge 3\), \(b, b', X' > 0\) with \(b'X' \le xb\), \(0 \le \Delta \le b\), \(\vartheta \ge 0\), and \(m \ge 1\).  Let \(\lvert a_t/b - x \rvert \le 1\) for \(t \in T\), and for \(u \in A_0\), \(1 \le j \le m\) let \(w_{u,j} \colon \Phi \to (0,\infty)\) satisfy \(b \le w_{u,j} \le b + \Delta\) on the thick cells for \(u\) and \(w_{u,j} \le 2b'\) on the others.  Suppose \(\lvert G_\theta(y)\rvert \le \vartheta y\) for all \(\theta\) and \(y \ge X'/3\), and put \(\mathcal{M} := \log(x+1) + C_d + \max \lvert c_{\chi_\theta}\rvert\) with \(C_d\) absolute, the maximum over the \(4o\) characters \(\theta_{\xi,j,l}\) with \(\xi \in \langle\psi\rangle\) (as in \(\mathcal{K}(x)\) of Problem~\ref{gf-pr:move-x}).  If the mean over \(u \in A_0\) of \(\sum_{t,j}\Re\De(a_t, t, w_{u,j})\) is at most \(\mathrm{Err}\), then
\begin{multline*}
  N_T m b\,\frac{\delta}{2}\Bigl(\sum_{\xi\in\langle\psi\rangle}\kappa_\xi\Re G_\xi(x) - \nu\sum_{\xi\in\langle\psi\rangle}\lvert G_\xi(x)\rvert - \sum{}' \lvert G_\theta(x)\rvert - 4o\,\mathcal{M}\Bigr) \\
  \le \mathrm{Err} + N_T m\bigl(4n\vartheta(x+1)b + 12n(x+1)\Delta\bigr),
\end{multline*}
where \(\sum'\) runs over the \(3o\) characters \(\theta_{\xi,j,l}\), \(\xi \in \langle\psi\rangle\), \((j,l) \ne (0,0)\).
\end{theorem}

\noindent\emph{Simplest incarnation}: Theorem~\ref{gf-thm:two-scales}.  \emph{Status}: proved in \cite{Adr}.  Here \(\vartheta\), \(\Delta\), \(\mathcal{M}\) are the numbers \(\omega\), \(\eta'\), \(M\) of \cite{Adr}.  Adriazola then divides by \(N_Tmb\delta/2\): if the sizes satisfy \cite[(4.6), (4.7)]{Adr} and \(\mathrm{Err} \le x^{C\varepsilon}\Bc\), where
\[
  \Bc := N_T a\Bigl(1 + \frac{m^2}{k}\Bigr) + m\Bigl(2nmb\,x^\varepsilon + \frac{\delta n\,k\,b\,x^{2\varepsilon}}{n^2}\Bigr) ,
\]
and if \(4n(x+1)\vartheta \le 1/x\), \(12n(x+1)\Delta \le b/x\) and \(4x^{\varepsilon-1} + \nu\sum_\xi\lvert G_\xi(x)\rvert + 4o\mathcal{M} \le \frac3{13}x^{1/2+(C+3)\varepsilon}\), then \(\sum_\xi\kappa_\xi\Re G_\xi(x) \le x^{1/2+(C+3)\varepsilon} + \sum'\lvert G_\theta(x)\rvert\) \cite[Prop.~4.3 and Cor.~4.4]{Adr}.  Status: proved in \cite{Adr}.

\begin{problem}\label{pr:excess-avg}
(i) For \(\psi = 1\), \(\De(a,t,w) = \sum_{\rho\in\Phi} w(\rho)R_{t,\rho}(a/w(\rho))\); so Theorem~\ref{thm:excess} is Theorem~\ref{gf-thm:balance}.
(ii) Take \(\psi = 1\), so that \(A_0 = H\), and \(m = 1\), \(T = \{t\}\), thick weights \(b\).  The mean over \(u \in H\) of the thick part of \(\De(a, t, w_u)\) is \(\frac{\delta b}{2}\sum_{(j,l)}\theta_{1,j,l}(\iota_d)\,G_{\theta_{1,j,l}}(a/b) = \frac{\delta b}{2} R_{q,d}(a/b)\), which is Problem~\ref{gf-pr:thick-avg}.
\end{problem}

The set \(T\) is chosen with prescribed shares on the cosets of \(A_0\).

\begin{theorem}[{\cite[Lemma~5.18]{Adr}}]\label{thm:shares}
If for every \(\gamma \in H\) the number of \(t \in T\) in \(\gamma A_0\) differs from \(\lambda_\gamma N_T\) by at most \(\delta' N_T\), where \(\lambda_\gamma = 4/(3o+1)\) for \(\gamma \in A_0\) and \(3/(3o+1)\) otherwise, then \(\lvert\sum_{t\in T}\xi(t) - N_T/(3o+1)\rvert \le o\,\delta' N_T\) for \(\xi \in \langle\psi\rangle \smallsetminus 1\).
\end{theorem}

\noindent\emph{Simplest incarnation}: Problem~\ref{gf-pr:shares}, with \(\delta' = 0\).  \emph{Status}: proved in \cite{Adr}; elementary.

\begin{problem}\label{pr:shares-nu}
With \(\delta' = x^{-1+\varepsilon}\), Theorem~\ref{thm:shares} gives the hypothesis of Theorem~\ref{thm:main-step} with \(\kappa = 1/(3o+1)\) and \(\nu = o\,x^{-1+\varepsilon}\).  If \(\lvert G_\xi(y)\rvert \le B_\psi\, y\) for \(\xi \in \langle\psi\rangle\), then \(\nu\sum_\xi\lvert G_\xi(x)\rvert \le o^2B_\psi\,x^\varepsilon\).
\end{problem}

So the \(\nu\) term costs \(x^\varepsilon\), far below \(x^{1/2}\).  What remains are the \(3o\) values \(\lvert G_\theta(x)\rvert\) in \(\sum'\), at the scale \(x\), and the thin-scale bound, at scales \(y \ge X'/3\).

\section{The auxiliary characters}\label{sec:aux}

The characters in \(\sum'\) have conductors divisible by \(q\) or \(d\), larger than any power of \(x\); no zero-free region for a fixed \(L\)-function bounds them at one point.  The primes \(q, d\) are therefore chosen, as in Chapters~\ref{zi-chap} and~\ref{rc-chap}, so that the bound holds on a whole window.  The classical inputs, with absolute constants, are listed in \cite[\S5.1]{Adr}: Chebyshev's bounds; the prime number theorem for characters with the exceptional zero \cite[Ch.~20]{Dav}; Page's theorem \cite{Pa}; Montgomery's density theorem \cite[Thm.~12.2]{Mon}; and the prime number theorem in progressions in Linnik's range, which \cite[(C7)]{Adr} derives from Gallagher's theorem \cite[Thm.~7]{Gal} with the exceptional zero kept in the main term.  All are classical. % check-numbers: ok (Davenport): Ch.~20

\begin{theorem}[thin scales, {\cite[Prop.~5.3]{Adr}}]\label{thm:thin}
There are absolute \(0 < c_1 \le \frac12\) and \(C_1\) such that the following holds.  Let \(\chi\) be primitive modulo \(N \ge 3\), \(\lambda \ge \max(1, \log N)\) and \(0 < c_\ast \le \frac12\), and suppose \(\chi\) is not real or \(L(\beta,\chi) \ne 0\) for \(1 - c_\ast < \beta < 1\).  Then for \(y \ge 2\) and \(t_0 := \log y\),
\[
  \frac{\lvert G_\chi(y)\rvert}{y} \le \frac{2e^{-c_\ast t_0}}{c_\ast} + C_1\lambda^4\Bigl(\frac{2\lambda}{c_1}e^{-c_1t_0/(2\lambda)} + \frac{8}{c_1^2}\Bigl(1 + \frac{c_1\sqrt{t_0}}{2}\Bigr)e^{-c_1\sqrt{t_0}/2}\Bigr) .
\]
The same decay holds for \(G\).
\end{theorem}

\noindent\emph{Simplest incarnation}: the thin-scale bound of Section~\ref{gf-sec:thick}, which Chapter~\ref{gf-chap} assumed.  \emph{Status}: classical; proved in \cite{Adr} from the prime number theorem for characters.

\begin{theorem}[first prime, {\cite[Thm.~5.10]{Adr}}]\label{thm:first}
Let \(J \ge 1\) and \(0 < \tau \le \frac12\).  For \(Y\) large in terms of \(N_0, J, \varepsilon, \tau\), put \(z := Y^{\varepsilon/5}\) and let \(\mathcal{P}\) be the set of primes \(q \equiv 3 \pmod 8\) in \((e^{z/2}, e^z]\).  For all \(q \in \mathcal{P}\) outside a set of at most \(Y^{-\tau}\lvert\mathcal{P}\rvert\) elements, every primitive \(\xi\) of conductor dividing \(N_0\) and \(\chi := \xi\chi_q\) satisfy:
(i) \(L(s,\chi) \ne 0\) for \(\Re s > 1 - \frac1{400}\), \(\lvert\Im s\rvert \le N_0e^z\);
(ii) \(\lvert c_\chi\rvert \le 401\,c_0\,J\log Y\), with \(c_0\) the Chebyshev constant;
(iii) \(\lvert G_\chi(y)\rvert \le y^{1/2+\tau}\) for every \(y \in [Y, Y^J]\).
\end{theorem}

\noindent\emph{Simplest incarnation}: Theorem~\ref{rc-thm:estB-q}, and Theorem~\ref{zi-thm:good-q-interval}.  \emph{Status}: proved in \cite{Adr} from the classical inputs above; Gallagher's density theorem is not needed for \(q\).

\begin{theorem}[second prime, {\cite[Thm.~5.12 and Cor.~5.13]{Adr}}]\label{thm:second}
Let \(J, \tau, Y, z\) be as in Theorem~\ref{thm:first}, let \(q \equiv 3 \pmod 8\) be a prime in \((e^{z/2}, e^z]\), let \(z\log^2Y \le z' \le J^2z\log^2Y\), and let \(\mathcal{P}'\) be the set of primes \(d \in (e^{z'/2}, e^{z'}]\) with \(d \equiv 7 \pmod{16}\) and \(d \equiv -1\) modulo \(q\) and modulo every odd prime \(p \le z\).  For all \(d \in \mathcal{P}'\) outside a set of at most \(Y^{-\tau}\lvert\mathcal{P}'\rvert\) elements, conclusions (i)--(iii) of Theorem~\ref{thm:first} hold for \(\chi := \xi\chi_q^a\chi_d\), \(a \in \{0,1\}\), with \(N_0e^{z+z'}\) in (i).  Moreover, if (i) holds for \(\xi\chi_q\) and for \(\xi\chi_q^a\chi_d\), then every real primitive \(\chi\) whose conductor is divisible by \(q\) or \(d\) and divides \(2^jMqd\) has \(L(\beta,\chi) \ne 0\) for \(1 - \frac1{400} < \beta < 1\).
\end{theorem}

\noindent\emph{Simplest incarnation}: Section~\ref{rc-sec:linnik} and Theorem~\ref{rc-thm:linnik}.  \emph{Status}: proved in \cite{Adr} from the classical inputs, with Linnik's range through Gallagher's theorem; the constants are not effective, by Siegel's theorem.  The congruences on \(d\) fix \(\bigl(\frac sd\bigr)\) for every \(s\) whose prime factors are at most \(z\); this is why \(\log d\) must exceed \(\log\) of the modulus \(16q\prod_{p\le z}p\), of order \(z\), by the factor \(\log^2Y\).

\begin{problem}\label{pr:estB-base}
(i) For \(M = 1\), \(J = \frac32\) and \(\tau = \varepsilon\), the case \(\xi = 1\) of Theorem~\ref{thm:first}(iii) is the bound of Theorem~\ref{rc-thm:estB-q} for \(A = \frac32\), with \(q \equiv 3 \pmod 8\) and \(\log q \in (z/2, z]\) in place of \(\log q\) of size \((\log Y)^3\).
(ii) If the conclusions hold at a scale \(Y_0\) with exponent \(J\), and \(Y_0 \le Y\), \(Y^2 \le Y_0^J\), then they hold on the window \([Y, Y^{3/2}]\); with \(z' = z\log^2Y_0\) there, \(\log q \le Y_0^{\varepsilon/5}\) and \(\log d \le Y_0^{\varepsilon/5}\log^2Y_0\).
\end{problem}

\begin{problem}\label{pr:twists}
Take \(N_0 = 8M\).  Each of the \(3o\) characters in \(\sum'\) is \(\xi\chi_q^j\chi_d^l\) with \(\xi\) primitive inducing a power of \(\psi\) and \((j,l) \ne (0,0)\), and Theorems~\ref{thm:first}(iii) and~\ref{thm:second}(iii) bound all of them: \(\lvert G_\theta(x)\rvert \le x^{1/2+\varepsilon}\), so \(\sum'\lvert G_\theta(x)\rvert \le 3o\,x^{1/2+\varepsilon}\).  For \(\psi = 1\) these are \(\chi_q, \chi_d, \chi_q\chi_d\); for \(\psi = \chi_8\) there are six.
\end{problem}

\begin{problem}\label{pr:thin}
Let \(\log N \le \lambda := A\,x^{\varepsilon/5}\log^2x\) with \(A\) fixed, and \(t_0 \ge x^{\varepsilon/3} - \log 3\).  For \(x\) large in terms of \(A\), \(c_\ast\), \(\varepsilon\), the right side of Theorem~\ref{thm:thin} is at most \(\exp(-x^{\varepsilon/8})\).  The exponents that matter are \(\frac13 - \frac15 = \frac{2}{15} > \frac18\) and \(\frac16 > \frac18\).
\end{problem}

By Theorem~\ref{thm:second}, the real characters of conductor divisible by \(q\) or \(d\) have no zero in \((1 - \frac1{400}, 1)\), and the finitely many others depend on \(M\) alone; so \(c_\ast\) depends on \(M\) only.  With \(\log N \le \lambda\) (Problem~\ref{pr:rd}), Problem~\ref{pr:thin} gives the thin-scale bound with \(\vartheta = \exp(-x^{\varepsilon/8})\) for all \(4n\) characters \(\theta\) of \(\Gamma\) \cite[proof of Thm.~6.3]{Adr}.

\section{The host field and the data}\label{sec:data}

\begin{theorem}[host field, {\cite[Prop.~5.14]{Adr}}]\label{thm:field}
Let \(q \ne d\) be primes \(\equiv 3 \pmod 4\) not dividing \(M\), \(e \ge 0\), \(N_1 := \operatorname{lcm}(2^{e+2}, M)\), \(U_1 := \{a \in (\Z/N_1)^\times : a \equiv \pm1 \bmod 2^{e+2},\ \psi(a) = 1\}\), \(H := (\Z/N_1)^\times/U_1\), and let \(E \subset \Q(\zeta_{N_1qd})\) be the field fixed by the classes that lie in \(U_1\) modulo \(N_1\) and are squares modulo \(q\) and \(d\).  Then \(E\) is a CM field containing \(\sqrt{-q}, \sqrt{-d}\), with group \(H \times \{\pm1\}^2\); \(n := \lvert H\rvert = c\,2^e\) with \(c \mid o\); the image \(A_0\) of \(\ker\psi\) is cyclic of index \(o\); and the characters of \(H\) trivial on \(A_0\) are the powers of \(\psi\).
\end{theorem}

\noindent\emph{Simplest incarnation}: Problem~\ref{gf-pr:galois}, the case \(\psi = 1\), where \(H = (\Z/4n\Z)^\times/\{\pm1\}\).  \emph{Status}: classical (class field theory of \(\Q(\zeta_N)\)).

Let \(x\) be the scale, and put
\[
  b := e^{x^\varepsilon}, \qquad a := xb, \qquad X' := \exp(x^{\varepsilon/3}), \qquad b' := a/X', \qquad \delta_1 := x^{-\varepsilon/8} .
\]

\begin{theorem}[{\cite[Props.~5.15--5.17, 5.20 and~18.7, Cor.~18.8]{Adr}}]\label{thm:sizes}
For \(x\) large in terms of \(\varepsilon\) and \(\psi\), and \(n\) as in Theorem~\ref{thm:field} with \(x^{2+8\varepsilon} \le n \le 2x^{2+8\varepsilon}\):
(i) there are even \(m \in [\sqrt{x\delta_1}, 2\sqrt{x\delta_1}]\) and \(k \in [x^{9/2+10\varepsilon}, 2x^{9/2+10\varepsilon}]\), and with \(\lvert T\rvert \ge \frac12\delta_1n\), \(\delta \ge x^{-\varepsilon}\) they satisfy the size conditions \cite[(4.6), (4.7)]{Adr} of the division by \(N_Tmb\delta/2\);
(ii) there are nodes \(w_{u,j}\) with \(b \le w_{u,j} \le b + \Delta\) at the thick cells, \(1 \le w_{u,j} \le 2b'\) elsewhere, where \(\Delta := W + mh_0\), \(W\) is the number of roots of unity in \(E\) and \(h_0 := 2\lceil x^{402}\rceil\); and then \(12n(x+1)x\Delta \le b\) and \(4n(x+1)\exp(-x^{\varepsilon/8}) \le 1/x\);
(iii) for \(L \ge 2\) and \(n_0 \ge 1\) there is a set \(T_2'\) of integers \(s\) with \(4Ls < n_0\), meeting no line in three points as in Problem~\ref{cs-pr:no-relation}, and \(\lvert T_2'\rvert \ge (n_0/L)\exp(-36\sqrt{\log n_0\log L})\);
(iv) there is \(T \subset H\) with \(\frac12\delta_1 n \le \lvert T\rvert \le 2\delta_1 n\), satisfying Theorem~\ref{thm:shares} with \(\delta' = x^{-1+\varepsilon}\), and with few tops whose Frobenius orbits are not good \cite[Def.~5.19]{Adr}.
\end{theorem}

\noindent\emph{Simplest incarnations}: (iii) is Theorem~\ref{cs-thm:behrend}, (iv) is Section~\ref{cs-sec:tops}.  \emph{Status}: proved in \cite{Adr}; (iii) is Behrend's construction \cite{Beh}, (iv) a random choice.

\begin{problem}\label{pr:behrend}
Let \(n_0 = n/o\) with \(x^{2+8\varepsilon} \le n \le 2x^{2+8\varepsilon}\), and \(2 \le L \le C\log x\).  For \(x\) large in terms of \(\varepsilon\), \(o\), \(C\), the bound of Theorem~\ref{thm:sizes}(iii) is at least \(x^{-\varepsilon}n\); so \(\delta \ge x^{-\varepsilon}\) for \(T_2 := \{a_0^s : s \in T_2'\}\), \(a_0\) a generator of \(A_0\); the map \(s \mapsto a_0^s\) is injective on \(T_2'\) because \(0 \le s < n_0 = \lvert A_0\rvert\).
\end{problem}

\begin{problem}\label{pr:nodes}
Under Theorem~\ref{thm:sizes}(ii), with \(\vartheta = \exp(-x^{\varepsilon/8})\), \(X := X'/3\) and \(\delta \ge x^{-\varepsilon}\):
(i) \(a_t/w_{u,j}(\rho) \ge X\) at every thin cell, if \(\lvert a_t/b - x \rvert \le 1\) and \(x \ge 3\);
(ii) the last two terms of Theorem~\ref{gf-thm:two-scales} are at most \(3x^{\varepsilon - 1}\), using \(\sum_{r \le u}\Lambda(r) \le 1.04\,u\) \cite{RS}.
\end{problem}

\begin{problem}\label{pr:rd}
The root discriminant \(\rd_E := \lvert\operatorname{disc}E\rvert^{1/[E:\Q]}\) satisfies \(\log\rd_E \le \log(8Mn) + \log q + \log d\).  If \(\log q \le Y_0^{\varepsilon/5}\), \(\log d \le Y_0^{\varepsilon/5}\log^2Y_0\) and \(Y_0 \le x\), then \(\ell_a := \log(a\,\rd_E) \le 2x^\varepsilon\) for large \(x\) (the number \(\Lambda_a\) of \cite[Lemma~18.15(iii)]{Adr}), and every \(\chi_\theta\) has \(\log\) conductor at most \(4x^{\varepsilon/5}\log^2x\).
\end{problem}

Admissible data at the scale \(x\) for the window \([Y, Y^{3/2}]\) with base \(Y_0\) \cite[Def.~18.1]{Adr} consist of: \(q, d\) as in Theorems~\ref{thm:first} and~\ref{thm:second} at the scale \(Y_0\), with exponent \(J\) and \(z' = z\log^2Y_0\); \(e\) with \(x^{2+8\varepsilon} \le c\,2^e \le 2x^{2+8\varepsilon}\) and the field of Theorem~\ref{thm:field}; \(m, k, T_2, T\) as in Theorem~\ref{thm:sizes}; two level primes \(g_0, g_1\) of size \((kb)^8\) split in \(K\), from Thorner--Zaman \cite{TZ} with Stark's bound \cite{Sta}; three primes \(\equiv 1\) modulo a conductor, from Linnik's theorem \cite{Lin}; and condition (D6): the Tate--Shafarevich groups \(\mathrm{Sha}_l(E, i)\), \(1 \le i < x^{400}\), have total length \(\sum_l \log l\cdot\operatorname{length} \le \exp(\frac12x^\varepsilon)\).

\begin{theorem}[{\cite[Prop.~6.2 and Prop.~18.4]{Adr}}]\label{thm:data}
(i) For \(Y\) large and every \(x \in [Y, Y^{3/2}]\), the data other than (D6) exist with \(Y_0 = Y\).
(ii) Assume the Selmer bound on windows \cite[Hyp.~17.11]{Adr}.  There is \(J \ge 2\) such that for arbitrarily large \(Y\) there is \(Y_0 \le Y\) with \(Y^2 \le Y_0^J\) for which admissible data with base \(Y_0\) exist at every \(x \in [Y, Y^{3/2}]\).
\end{theorem}

\noindent\emph{Simplest incarnation}: the example \(q = 11\), \(d = 1319\), \(n = 4\) of Chapter~\ref{gf-chap}.  \emph{Status}: (i) is proved in \cite{Adr} from published theorems; (ii) is proved in \cite{Adr} from the Selmer bound on windows; in \cite{Adr} this is \cite[Hyp.~17.11]{Adr}, assumed there.  The formalization proves the final theorem; see Section~\ref{as-sec:map} of Chapter~\ref{as-chap}.  Condition (D6) is the only reason for the base \(Y_0\): Theorem~\ref{thm:first} must be applied at \(Y_0\), and Problem~\ref{pr:estB-base}(ii) carries it to the window.

\section{The last step}\label{sec:last}

Theorem~A of \cite{Adr} asserts: there is an absolute \(C\) such that for every \(\psi\) and small \(\varepsilon\) there is \(J\) with the following property, for arbitrarily large \(Y\) and some \(Y_0 \le Y\) with \(Y^2 \le Y_0^J\).  At every \(x \in [Y, Y^{3/2}]\) there are admissible data with base \(Y_0\), weights \(a \le a_t \le a + a/(kx^{13})\) and nodes as in Theorem~\ref{thm:sizes}(ii), with \(\sum_{t,j}\Re\De(a_t, t, w_{u,j}) \le x^{C\varepsilon}\Bc\) for every \(u \in A_0\).  One can take \(C = 3\).  \cite[\S18.6]{Adr} derives it from \cite[Hyp.~17.11]{Adr} and the chain records \cite[Hyp.~18.10]{Adr}, both assumed there.  The formalization proves the final theorem, with chain records only for data it chooses, a weaker statement than \cite[Hyp.~18.10]{Adr} (Section~\ref{as-sec:map} of Chapter~\ref{as-chap}).

\begin{theorem}[{\cite[Thm.~6.3]{Adr}}]\label{thm:last}
Theorem~A implies that \(L(s,\psi) \ne 0\) for \(\Re s > \frac12\).
\end{theorem}

\noindent\emph{Simplest incarnation}: Theorems~\ref{gf-thm:two-scales} and~\ref{zr-thm:dirichlet-intervals}.  \emph{Status}: proved in \cite{Adr} from the theorems of this chapter.

\begin{problem}\label{pr:theoremA}
Assume Theorem~A.  On its windows, for large \(Y\), \(G^\psi(x) \le x^{1/2+(C+4)\varepsilon}\).  The only input from Part~II is the bound for \(\sum\Re\De\); all other hypotheses of Theorem~\ref{thm:main-step} and of the division by \(N_Tmb\delta/2\) are supplied by Theorems~\ref{thm:shares}--\ref{thm:data} and Problems~\ref{pr:shares-nu}--\ref{pr:rd}.
\end{problem}

Since \(C\) does not depend on \(\varepsilon\), Theorem~\ref{thm:window} with \(\theta = \frac12 + (C+4)\varepsilon\) and \(\varepsilon \to 0\) gives Theorem~\ref{thm:last}.

\paragraph{The example \(\psi = \chi_8\).}  Take \(q = 11\), \(d = 1319\), \(n = 4\), \(H = \{[1],[3],[5],[7]\}\) and \(\psi := \xi^2\) as in Section~\ref{gf-sec:subgroup}, so that \(\chi_\psi = \chi_8\), \(o = 2\), \(A_0 = \{[1],[7]\}\).  No subset of \(H\) has the shares \(\frac47\), \(\frac37\) on the two cosets, so let \(T\) be the multiset \(\{[1],[1],[7],[7],[3],[3],[5]\}\), with sums and means over \(T\) counted with multiplicity; in the argument \(n\) is large and a set suffices.

\begin{problem}\label{pr:chi8-shares}
The mean of the thick part of~\eqref{gf-eq:balance} over \(u \in A_0\) and \(t \in T\) is \(\frac{\delta b}{2}R^\psi(a/b)\), with \(R^\psi := R_{q,d} + \frac17R_{\xi^2}\).  Its untwisted part is \(G + \frac17G_{\chi_8} = G^\psi\).
\end{problem}

\begin{problem}\label{pr:chi8-twists}
\(R^\psi - G^\psi = G_{\chi_{11}} - G_{\chi_{1319}} - G_\eta + \frac17(G_{\chi_8\chi_{11}} - G_{\chi_8\chi_{1319}} - G_{\chi_8\eta})\), with primitive characters of conductors \(11\), \(1319\), \(14509\), \(88\), \(10552\), \(116072\).  If these six satisfy \(\lvert G_\chi(x)\rvert \le x^{1/2+\varepsilon}\), then \(G^\psi(x) \le R^\psi(x) + \frac{24}{7}x^{1/2+\varepsilon}\).
\end{problem}

These six are the characters of Problem~\ref{pr:twists} for \(\psi = \chi_8\); here \(q\) and \(d\) are far too small, and the example only names them.

\begin{problem}\label{pr:chi8-positive}
\(L(\sigma, \chi_8) > 0\) for \(0 < \sigma < 1\).
\end{problem}

The argument does not use this: for a general \(\psi\) a real zero of \(L(s,\psi)\) is excluded only by the positive residue in Theorem~\ref{thm:window}.

\paragraph{Deliverable.}  This chapter proves none of the conditions (K1)--(K8) of a chain record \cite[Def.~18.9]{Adr}; it supplies what they are stated about and what follows from them.  It gives the admissible data, the weights and the nodes \(w_{u,j}\) (Theorems~\ref{thm:field}--\ref{thm:data}), with the numbers \(b\), \(b'\), \(a\), \(\Bc\) and \(\ell_a \le 2x^\varepsilon\) that enter (K1)--(K8); the sum \(S\) of (K1), which is the left side of Theorem~\ref{thm:excess} summed over \(t\) and \(j\); and the nodes in the first inequality of (K7).  Conversely, given chain records at every scale of a window, \cite[Props.~18.14 and~18.16]{Adr} give \(\sum\Re\De \le x^{3\varepsilon}\Bc\), and this chapter turns that into the Riemann hypothesis for \(L(s,\psi)\): Theorem~\ref{thm:sizes}(ii) with Problem~\ref{pr:nodes}(i) gives the normalization condition (V2), and Problem~\ref{pr:nodes}(ii) gives (V3), of the final chapter, \emph{Assembly at one scale}; Problems~\ref{pr:twists} and~\ref{pr:thin} give its window conditions (W1) and (W2); and Problem~\ref{pr:behrend} with Theorem~\ref{thm:sizes}(iv) gives \(\delta \ge x^{-\varepsilon}\) and \(N_T \ge \frac12x^{-\varepsilon/8}n\).  The only input here without a published proof is (D6): in \cite{Adr} it rests on \cite[Hyp.~17.11]{Adr}, assumed there; see Section~\ref{as-sec:map} of Chapter~\ref{as-chap}.

\section*{Hints}

\subsection*{\S\ref{sec:windows}}
\begin{description}
\item[\ref{pr:windows}] (i) \(\lvert G(x)\rvert \le Bx\) by Theorem~\ref{zr-thm:explicit}, since \(\lvert x^\rho\rvert \le x\); then Problem~\ref{zr-pr:F-poles}, and \(\rho(\rho-1) < 0\) for real \(\rho \in (0,1)\).  (ii) Problem~\ref{zr-pr:product}.
\end{description}

\subsection*{\S\ref{sec:average}}
\begin{description}
\item[\ref{pr:excess-avg}] (i) Put \(\rho := tgc\).  (ii) Orthogonality over \(u \in H\) leaves the characters trivial on \(H\); the thick cells for \(u\) are \(g\) with \(tgc\) in row \(1\), column in \(uT_2\).
\item[\ref{pr:shares-nu}] \(\lvert\langle\psi\rangle\rvert = o\).
\end{description}

\subsection*{\S\ref{sec:aux}}
\begin{description}
\item[\ref{pr:estB-base}] (ii) \(Y^{3/2} \le Y^2 \le Y_0^J\).
\item[\ref{pr:twists}] Problem~\ref{gf-pr:A0}; \(\chi_q\) and \(\chi_d\) are induced by Legendre symbols since \(q \equiv d \equiv 3 \pmod 4\).
\item[\ref{pr:thin}] Every prefactor is a power of \(x\).
\end{description}

\subsection*{\S\ref{sec:data}}
\begin{description}
\item[\ref{pr:behrend}] \(\sqrt{\log n\,\log\log x} = o(\log x)\).
\item[\ref{pr:nodes}] (i) \(b' = xb/X'\).  (ii) Use \(\Delta \le b/(12nx(x+1))\) and \(4n(x+1)\vartheta \le 1/x\).
\item[\ref{pr:rd}] Every character of \(E\) has conductor dividing \(N_1qd\), and \(2^e \le n\).
\end{description}

\subsection*{\S\ref{sec:last}}
\begin{description}
\item[\ref{pr:theoremA}] Divide by \(N_Tmb\delta/2\); bound \(\mathcal{M}\) by Theorems~\ref{thm:first}(ii) and~\ref{thm:second}(ii) for \((j,l) \ne (0,0)\), while the \(c_{\psi_j}\) depend on \(\psi\) alone, and \(\sum'\) by Problem~\ref{pr:twists}.
\item[\ref{pr:chi8-shares}] Problems~\ref{gf-pr:thick-A0} and~\ref{gf-pr:shares}.
\item[\ref{pr:chi8-twists}] Problem~\ref{gf-pr:A0}, and the triangle inequality.
\item[\ref{pr:chi8-positive}] Group the terms in blocks of eight, and use that \(t \mapsto t^{-\sigma}\) is convex.
\end{description}

\begin{thebibliography}{Gal}
\bibitem[Adr]{Adr} J.~Adriazola, \emph{The Alp\"oge proof of the Riemann hypothesis}, manuscript, revised 24 November 2026, \url{https://github.com/jimjamzola/RH-Alpoge-Reconstruction}.
\bibitem[Beh]{Beh} F.~A.~Behrend, \emph{On sets of integers which contain no three terms in arithmetical progression}, Proc.\ Nat.\ Acad.\ Sci.\ U.S.A.\ \textbf{32} (1946), 331--332.
\bibitem[Dav]{Dav} H.~Davenport, \emph{Multiplicative Number Theory}, 3rd ed., Grad.\ Texts in Math.\ 74, Springer, 2000.
\bibitem[Gal]{Gal} P.~X.~Gallagher, \emph{A large sieve density estimate near \(\sigma = 1\)}, Invent.\ Math.\ \textbf{11} (1970), 329--339.
\bibitem[Lean]{Lean} The Lean community, \emph{riemann-hypothesis}, a Lean formalization of Alp\"oge's argument, \url{https://github.com/leanprover-community/riemann-hypothesis}.
\bibitem[Lin]{Lin} Yu.~V.~Linnik, \emph{On the least prime in an arithmetic progression.~I.~The basic theorem}, Rec.\ Math.\ [Mat.\ Sbornik] N.S.\ \textbf{15} (1944), 139--178.
\bibitem[Mon]{Mon} H.~L.~Montgomery, \emph{Topics in Multiplicative Number Theory}, Lecture Notes in Math.\ 227, Springer, 1971.
\bibitem[Pa]{Pa} A.~Page, \emph{On the number of primes in an arithmetic progression}, Proc.\ London Math.\ Soc.\ (2) \textbf{39} (1935), 116--141.
\bibitem[RS]{RS} J.~B.~Rosser and L.~Schoenfeld, \emph{Approximate formulas for some functions of prime numbers}, Illinois J.\ Math.\ \textbf{6} (1962), 64--94.
\bibitem[Sta]{Sta} H.~M.~Stark, \emph{Some effective cases of the Brauer--Siegel theorem}, Invent.\ Math.\ \textbf{23} (1974), 135--152.
\bibitem[TZ]{TZ} J.~Thorner and A.~Zaman, \emph{A unified and improved Chebotarev density theorem}, Algebra Number Theory \textbf{13} (2019), 1039--1068.
\bibitem[Wa]{Wa} L.~C.~Washington, \emph{Introduction to Cyclotomic Fields}, 2nd ed., Grad.\ Texts in Math.\ 83, Springer, 1997.
\end{thebibliography}

\end{document}
